% Generated draft: requires mathematical and source review.
\documentclass[11pt]{article}
\usepackage[margin=25mm]{geometry}
\usepackage{fontspec}
\setmainfont{Latin Modern Roman}
\usepackage{amsmath,amssymb,mathtools,unicode-math}
\setmathfont{Latin Modern Math}
\usepackage{enumitem,needspace,xurl,hyperref}
\setlength{\parindent}{0pt}
\setlength{\parskip}{5pt}
\setlength{\emergencystretch}{3em}
\newcommand{\N}{\mathbb{N}}
\newcommand{\Z}{\mathbb{Z}}
\newcommand{\Q}{\mathbb{Q}}
\newcommand{\R}{\mathbb{R}}
\newcommand{\C}{\mathbb{C}}
\newcommand{\F}{\mathbb{F}}
\newcommand{\A}{\mathbb{A}}
\newcommand{\PP}{\mathbb{P}}
\newcommand{\E}{\mathbb{E}}
\newcommand{\eps}{\varepsilon}
\newcommand{\dd}{\,\mathrm{d}}
\begin{document}
\section*{UnsolvedMath ID: 30004164}
\textbf{Source problem code:} catalogue entry
\subsection{Definitions and mathematical statement}
Under the customary commutative, unital interpretation, an algebra $A$ essentially of finite type over a field $k$ is a localization $S^{-1}B$ of a finitely generated commutative $k$-algebra $B$. Such an algebra is Noetherian. Its Krull dimension $\dim A$ is the supremum of lengths of strictly increasing chains of prime ideals. For an integer $n\geq 1$, call $A$ \emph{$K_n$-regular} if, for every integer $m\geq 1$, the inclusion into the polynomial algebra induces an isomorphism of algebraic $K$-groups
\[
K_n(A)\xrightarrow{\;\cong\;}K_n\bigl(A[t_1,\ldots,t_m]\bigr).
\]
Call $A$ \emph{regular} if, for every prime ideal $\mathfrak p\subseteq A$, the local ring $A_{\mathfrak p}$ is regular local; equivalently, its maximal ideal has a minimal generating set of cardinality $\dim A_{\mathfrak p}$.

\textbf{Question (supplied formulation, associated in the background with Vorst's conjecture).} For every field $k$ and every algebra $A$ essentially of finite type over $k$, with finite Krull dimension $d\geq 0$, does the implication
\[
A\text{ is }K_{d+1}\text{-regular}
\quad\Longrightarrow\quad
A\text{ is regular}
\]
hold? No restriction on the characteristic of $k$ is stated in the supplied formulation.

\textbf{Unresolved source qualifications.} The commutative, unital conventions and the definitions above give the customary interpretation of the question; these conventions are not explicit in the supplied statement. The zero-ring convention is unspecified. More substantially, the supplied background says that the extracted report passage ends before the conjecture is stated, whereas the registry supplies a complete question and labels its statement as corrected and verified. Consequently, the question above records the supplied mathematical target, but its exact correspondence with the report remains unresolved. The supplied open label is a curation classification, not an independently verified mathematical status.
\subsection{Short English statement}
The supplied question asks whether invariance of $K_{d+1}$ under adjoining any finite number of polynomial variables forces a $d$-dimensional algebra essentially of finite type over a field to be regular. This uses the customary commutative, unital interpretation. The background reports a truncated source passage, so the exact report formulation and the supplied open status require independent review.
\subsection{Sources}
\begin{itemize}
\item Supplied local UnsolvedMath registry record, permanent ID $30004164$, source problem code \texttt{catalogue entry}, title \emph{K-Regularity Criterion for Ring Regularity}, in the Oberwolfach Reports Open Problems collection. The record lists proposed year $2019$ and supplies no named proposer. Its background associates the question with Vorst's conjecture.
\item Supplied report citation: \emph{Algebraic K-theory (2020)}, Oberwolfach Reports, pp.~1737--1790, DOI \texttt{10.4171/owr/2019/29}. Source URL and supplied literature-source link labeled ``OWR primary report'': \url{https://doi.org/10.4171/owr/2019/29}. This reference is supplied provenance; it has not been visited or independently verified here.
\item Supplied literature assessment dated \texttt{2026-08-22}: the extracted passage introduces Vorst's conjecture but is truncated before stating it; the record classifies the item as open under its queue rule because the exact claim is unclear. This is supplied dated curation, not an independently verified assessment of the literature or mathematical status.
\item The registry labels the statement \texttt{corrected\_verified} and reports manual semantic, typographic, and LaTeX review. That supplied review claim does not resolve the background's qualification concerning the truncated report passage.
\end{itemize}
\end{document}
